\chapter{鲁棒性与敏感性分析}

\section{参数敏感性定义}
对参数 $\theta_j$ 的局部敏感性定义为
\[
S_j(t)=\frac{\partial x(t)}{\partial \theta_j}\cdot \frac{\theta_j}{x(t)+\epsilon},
\]
其中 $\epsilon$ 为数值稳定项。

\section{噪声传播分析}
设观测噪声协方差为 $\mathbf{\Sigma}_y$，在线性近似下参数协方差可估计为
\[
\mathbf{\Sigma}_\theta \approx \left(\mathbf{J}^\mathsf{T}\mathbf{\Sigma}_y^{-1}\mathbf{J}\right)^{-1}.
\]
据此可评估参数置信区间与相关性。

\section{初值扰动试验}
在基准初值附近施加随机扰动：
\[
\boldsymbol{\theta}_0^{(j)}=\boldsymbol{\theta}_0+\delta^{(j)},\quad \delta^{(j)}\sim \mathcal{N}(0,\sigma_0^2\mathbf{I}),
\]
记录多次求解结果统计量，以判断算法对初始化的敏感程度。

\section{工程建议}
\begin{itemize}
  \item 使用参数归一化降低尺度差异带来的病态性；
  \item 优先选取信息量高的观测时段提升可辨识性；
  \item 对高噪声数据采用分段拟合与正则化联合策略。
\end{itemize}
